\chapter{Banks, Runs and Financial Crises}\label{ch:b2:banks-and-crises}

Banks borrow from depositors who can withdraw at any time and lend for years. That arrangement is useful, and it is
fragile: if enough depositors withdraw at once, a bank that would have paid everyone in full can fail. This chapter builds
the model of that fragility, looks at the private and public answers to it, and measures what banking crises have cost
over a century and a half. The chart counts them.

\begin{openingchart}{Number of systemic banking crises starting in each decade, 18 advanced economies, 1870s to
  2010s.}{Jord\`a--Schularick--Taylor Macrohistory Database (crisisJST); our computation}
\begin{tikzpicture}
\begin{axis}[omrecord, width=0.72\linewidth, height=0.4\linewidth, ybar, bar width=9pt, xmin=1862, xmax=2018,
    ymin=0, ymax=16, xtick={1870,1890,...,2010}, xticklabel style={/pgf/number format/1000 sep=},
    xticklabel={\pgfmathprintnumber{\tick}s}, ylabel={Crises}, nodes near coords,
    nodes near coords style={font=\scriptsize}]
  \addplot[fill=omDef!70, draw=omDef] table[x=decade, y=crises] {figdata/bachelor-2/29-banks-and-crises/out/decades.csv};
\end{axis}
\end{tikzpicture}
\end{openingchart}

% figdata: decades (total 88; 1920s 13, 1930s 14, 2000s 13; 1940s-1960s 0)
The eighteen economies had 88 systemic banking crises between 1870 and 2020: 13 in the 1920s, 14 in the 1930s, 13 in
the 2000s, and none at all from 1940 to 1969. Crises are not rare accidents of a distant past, and their frequency has
changed with the way banks were run and regulated.

\begin{recall}
From the school volumes: a bank's balance sheet, reserves, \omterm{def:b2:banks-and-crises:bank-run}{bank runs}, \omterm{def:b2:banks-and-crises:deposit-insurance}{deposit insurance} and the \omterm{def:b2:banks-and-crises:llr}{lender of last resort}, told
in words. From \cref{ch:b2:expected-utility}: \omterm{def:b2:expected-utility:risk-aversion}{risk aversion} and \omterm{def:b2:expected-utility:crra}{constant relative risk aversion}. From
\cref{ch:b2:three-equation-model}: the policy rate.
\end{recall}

\section{Short deposits, long loans}\label{sec:b2:banks-and-crises:maturity}

\begin{definition}[Maturity transformation]\label{def:b2:banks-and-crises:maturity-transformation}
\emph{Maturity transformation}\index{maturity transformation} is the financing of long-term, illiquid assets with
short-term liabilities that can be withdrawn on demand or at short notice.
\end{definition}

A household does not know whether it will need its savings next month or in ten years. Holding cash is safe but earns
nothing; lending to a firm for ten years earns more but cannot be undone cheaply. A bank pools many households: since only
some of them need their money at any date, it can keep a small part of its assets liquid and lend the rest long. Each
depositor gets a claim that can be cashed at any time and still earns part of the return on long loans. The same pooling is
the source of fragility: the bank can meet the withdrawals it expects, not withdrawals by everyone.

\section{The bank-run model}\label{sec:b2:banks-and-crises:dd}

\begin{definition}[Diamond--Dybvig model]\label{def:b2:banks-and-crises:dd}
The \emph{Diamond--Dybvig model}\index{Diamond--Dybvig model} is a three-date model of a bank that offers demand deposits
to depositors who learn privately whether they must consume early, and that invests in a technology that pays more if held
to the end than if liquidated early.
\end{definition}

There are three dates, 0, 1 and 2, and a large number of depositors, each with one unit at date 0. At date 1 each learns
privately whether he is \emph{impatient}, needing to consume at date 1 (a share $\pi$ of them), or \emph{patient}, able to
wait until date 2. A unit invested at date 0 returns $R > 1$ at date 2, or 1 if it is liquidated at date 1. Utility is
$u(c) = c^{1-\gamma}/(1 - \gamma)$, with $\gamma$ the \omterm{def:b2:expected-utility:rra}{coefficient of relative risk aversion}.

\begin{omfigure}
\begin{tikzpicture}[econ, every node/.style={font=\small}]
  \draw[thick, ->] (0,0) -- (11,0);
  \foreach \x/\t in {0.5/0, 5/1, 9.8/2} {\draw[thick] (\x,0.15) -- (\x,-0.15); \node[below] at (\x,-0.2) {date \t};}
  \node[above, align=center, text width=3.2cm] at (0.5,0.3) {each depositor\\ deposits 1};
  \node[above, align=center, text width=4.2cm] at (5,0.3) {types learnt privately;\\ withdrawers paid $c_1$\\ (liquidation yields 1)};
  \node[above, align=center, text width=3.4cm] at (9.8,0.3) {assets yield $R$;\\ the rest is shared\\ among those who waited};
\end{tikzpicture}

{\small The timing of the model. A deposit is a promise of $c_1$ to anyone who withdraws at date 1; those who wait share
what is left at date 2.}
\end{omfigure}

Without a bank each depositor invests alone: an impatient one liquidates and gets 1, a patient one waits and gets $R$. A
risk-averse depositor would rather narrow that gap: give up some of $R$ in the good case to get more than 1 in the bad
one. With $\pi$ known, a bank can do it. If it pays $c_1$ to the impatient and $c_2$ to the patient, it must liquidate
$\pi c_1$ at date 1 and keep $1 - \pi c_1$ invested, so that $(1 - \pi)c_2 = R(1 - \pi c_1)$, or
\begin{equation}\label{eq:b2:banks-and-crises:budget}
\pi c_1 + (1 - \pi)\frac{c_2}{R} = 1.
\end{equation}
Maximising $\pi u(c_1) + (1 - \pi)u(c_2)$ under \eqref{eq:b2:banks-and-crises:budget} gives $u'(c_1) = Ru'(c_2)$, so
$c_2 = R^{1/\gamma}c_1$ and
\begin{equation}\label{eq:b2:banks-and-crises:contract}
c_1 = \frac{1}{\pi + (1 - \pi)R^{1/\gamma - 1}}.
\end{equation}
If $\gamma > 1$, $R^{1/\gamma - 1} < 1$ and $c_1 > 1$: the bank pays early withdrawers more than the liquidation value of
their deposit. That is the insurance, and that is the weak point.

\begin{definition}[Bank run]\label{def:b2:banks-and-crises:bank-run}
A \emph{bank run}\index{bank run} is an equilibrium in which depositors who do not need their funds withdraw them because
they expect others to withdraw, so that withdrawals exceed what the bank can pay at the promised rate.
\end{definition}

Suppose a share $f \geq \pi$ of depositors withdraws at date 1. The bank liquidates $fc_1$ and, if $fc_1 < 1$, those who
waited share $R(1 - fc_1)$, each getting
\[
w(f) = \frac{R(1 - fc_1)}{1 - f}.
\]
If $f = \pi$, $w = c_2 > c_1$: a patient depositor who expects only the impatient to withdraw does best to wait. But $w$
falls as $f$ rises whenever $c_1 > 1$, and is zero once $f \geq 1/c_1$, when the bank has liquidated everything. A patient
depositor who expects many others to withdraw does best to withdraw too. Waiting is the \omterm{def:b2:static-games:best-response}{best response} if and only if
$w(f) \geq c_1$, that is
\begin{equation}\label{eq:b2:banks-and-crises:threshold}
f \leq f^* = \frac{R - c_1}{c_1(R - 1)}.
\end{equation}
Both outcomes are equilibria: everyone patient waits, and $f = \pi$; or everyone withdraws, and the bank fails. In the run
equilibrium the bank liquidates investments that would have returned $R$: the run destroys real output, not merely
transfers it.

\omcode{code/bachelor-2/29-banks-and-crises/listings/dd.py}{1}{32}{The optimal deposit contract and the run threshold, for
three degrees of \omterm{def:b2:expected-utility:risk-aversion}{risk aversion}.}\label{lst:b2:banks-and-crises:dd}

\omoutput{code/bachelor-2/29-banks-and-crises/outputs/dd.txt}{With $\pi = 0.25$ and $R = 2$: the more risk-averse the
depositors, the more the contract insures them and the lower the share of withdrawals that makes waiting worse than
running.}\label{out:b2:banks-and-crises:dd}

% code: listings/dd.py
With log utility ($\gamma = 1$) the contract is $c_1 = 1$, $c_2 = R$: there is no insurance, and no run, since waiting
pays $R$ whatever others do. With $\gamma = 2$ the bank promises 1.282 at date 1 and 1.812 at date 2; waiting stops paying
once 56.1 per cent of depositors withdraw, and the bank is empty at 78.0 per cent. With $\gamma = 4$ the threshold falls
to 39.2 per cent. The more insurance the contract gives, the more fragile the bank.

\begin{omfigure}
\begin{tikzpicture}[econ]
\begin{axis}[omrecord, grid=none, xmin=0, xmax=1, ymin=0, ymax=2.2, width=0.6\linewidth, height=0.4\linewidth,
    xlabel={Share of depositors withdrawing at date 1, $f$}, ylabel={Payoff}, clip=false,
    legend style={at={(1.03,0.5)}, anchor=west, font=\small}]
  \addplot[omDef, thick, domain=0:0.7803, samples=100] {2*(1 - 1.2815*x)/(1 - x)};
  \addlegendentry{wait, $w(f)$}
  \addplot[omThm, thick, dashed, domain=0:1, samples=2] {1.2815};
  \addlegendentry{withdraw, $c_1$}
  \addplot[omDef, thick, forget plot, domain=0.7803:1, samples=2] {0};
  \node[eqpoint] at (axis cs:0.25,1.8123) {};
  \node[above right, font=\small] at (axis cs:0.25,1.8123) {$f = \pi$: $c_2$};
  \node[eqpoint] at (axis cs:0.5607,1.2815) {};
  \draw[guide] (axis cs:0.5607,1.2815) -- (axis cs:0.5607,0);
  \node[left, font=\small] at (axis cs:0.5607,0.2) {$f^* = 0.561$};
  \draw[guide] (axis cs:0.7803,0) -- (axis cs:0.7803,0.25) node[above, font=\small] {$1/c_1$};
\end{axis}
\end{tikzpicture}

{\small Payoffs to a patient depositor, $\pi = 0.25$, $R = 2$, $\gamma = 2$. Left of $f^* = 0.561$ waiting pays more than
withdrawing; right of it, withdrawing pays more. Both ends are equilibria.}
\end{omfigure}

\section{Deposit insurance and the lender of last resort}\label{sec:b2:banks-and-crises:insurance}

The run equilibrium exists because a patient depositor fears that the bank will have nothing left at date 2. Remove the
fear and the run disappears. Three answers have been used.

\emph{Suspension of convertibility.} The bank announces that it will pay at most $\pi c_1$ at date 1, then stop. Since
patient depositors know that the bank will keep $1 - \pi c_1$ invested, waiting pays at least $c_2$ whatever others do, and
none of them runs. It works when $\pi$ is known; if the share of impatient depositors is itself uncertain, some impatient
depositors who come after the stop go unpaid.

\begin{definition}[Deposit insurance]\label{def:b2:banks-and-crises:deposit-insurance}
\emph{Deposit insurance}\index{deposit insurance} is a guarantee, by a fund or the government, that deposits up to a limit
will be repaid in full whatever happens to the bank, financed by premiums paid by banks or by taxes.
\end{definition}

If depositors who wait are promised $c_2$, waiting is always the \omterm{def:b2:static-games:best-response}{best response} for a patient depositor, so no run starts
and the guarantee is never called upon. In the model insurance costs nothing in equilibrium.

\begin{definition}[Lender of last resort]\label{def:b2:banks-and-crises:llr}
A \emph{lender of last resort}\index{lender of last resort} is an institution, usually the central bank, that lends to
solvent banks facing withdrawals they cannot meet from liquid assets, against good collateral, when markets will not.
\end{definition}

A \omterm{def:b2:banks-and-crises:llr}{lender of last resort} lets the bank pay early withdrawers without liquidating its investments at a loss, so the payoff
to waiting no longer falls as $f$ rises. Its condition is solvency: the bank's assets must be worth enough at date 2 to
repay the loan. In practice, telling an illiquid bank from an insolvent one in the middle of a run is hard.

\begin{model}{Diamond--Dybvig model}
\assumes Depositors with private liquidity needs; an investment that pays more if held to the end; demand deposits paid in
the order of arrival; the share of impatient depositors known.
\predicts A deposit contract that insures depositors; a second equilibrium, the run, in which a sound bank fails; no run
under suspension of convertibility or \omterm{def:b2:banks-and-crises:deposit-insurance}{deposit insurance}.
\tested About 1,350 American banks suspended in 1930, 2,300 in 1931, 1,450 in 1932 and 4,000 in 1933, 3,800 of them by 16
March; federal \omterm{def:b2:banks-and-crises:deposit-insurance}{deposit insurance} took effect on 1 January 1934, and the insurer handled nine failures of insured banks
that year. The reforms of 1933--34 came together, so this is consistent with the prediction, not a clean test of it.
\breaks The model does not say which equilibrium occurs, so on its own it cannot predict when a run will happen. Insurance
that removes the run also removes depositors' reasons to watch what their bank does with their money, which invites
risk-taking (the debate below). And the bank's assets are safe in the model; in the record the doubt that starts a run is
usually about the assets.
\end{model}

% ledger: b2:dd1983, b2:fdic1930s
Diamond and Dybvig (1983) showed that deposit contracts can give depositors better \omterm{def:b2:exchange-economies:allocation}{allocations} than markets for the
underlying assets, that these contracts have a run as a second equilibrium, that runs cause real damage, and that
government \omterm{def:b2:banks-and-crises:deposit-insurance}{deposit insurance} can, in some circumstances, do better than any private contract.

\begin{record}[Clearing houses in the panic of 1907]
% ledger: b2:panic1907, b2:gorton1985
Before 1914 the United States had no central bank. In October 1907 runs hit banks linked to a failed attempt to corner the
stock of a copper company; from 18 October the depositors of the Knickerbocker Trust withdrew nearly 8 million dollars in
four days, and it suspended on 22 October. The New York Clearing House, an association of banks set up to settle cheques
between members, declared the affected member banks solvent, forced out their management, and lent to them through clearing
house loan certificates; on 26 October its banks restricted the convertibility of deposits into cash. Industrial output fell
17 per cent in 1908 and real GNP 12 per cent. Gorton (1985) argued that in panics the clearing house turned its members
into something like a single firm producing \omterm{def:b2:banks-and-crises:deposit-insurance}{deposit insurance}: a private answer to the problem the model describes. The
panic was followed by the reform that created a public one, the Federal Reserve.
\end{record}

\begin{debate}[Does deposit insurance make crises more likely?]
% ledger: b2:dkd2002, b2:adz2014
Demirg\"u\c{c}-Kunt and Detragiache (2002), with data on 61 countries in 1980--97, found that explicit \omterm{def:b2:banks-and-crises:deposit-insurance}{deposit insurance}
tended to make banking crises more likely, the more so where bank interest rates were deregulated, institutions weak,
coverage extensive and the scheme run by the government. Anginer, Demirg\"u\c{c}-Kunt and Zhu (2014) found that more
generous insurance was associated with higher bank risk before 2007 and with lower risk and greater stability during the
crisis, with a destabilising net effect over the whole period, and that good supervision reduced the effect in good years.
Both rest on comparisons across countries whose insurance schemes were not assigned at random, and the size of the effect
is not settled.
\end{debate}

\section{Leverage and forced sales}\label{sec:b2:banks-and-crises:leverage}

\begin{definition}[Leverage ratio]\label{def:b2:banks-and-crises:leverage}
A bank's \emph{leverage ratio}\index{leverage ratio (bank)} is its equity divided by its total assets; its inverse, assets
per unit of equity, is called its leverage.
\end{definition}

A bank with assets of 100 financed by 96 of deposits and other debt and 4 of equity has a \omterm{def:b2:banks-and-crises:leverage}{leverage ratio} of 4 per cent and
a leverage of 25. A fall of 2 per cent in the value of its assets wipes out half its equity; a fall of 4 per cent wipes out
all of it. With 10 of equity the same 2 per cent fall costs a fifth. Equity is the cushion that keeps an asset loss from
becoming a loss to depositors, and so a reason for depositors to run.

\begin{definition}[Fire sale]\label{def:b2:banks-and-crises:fire-sale}
A \emph{fire sale}\index{fire sale} is a sale of assets at a price below their value in normal use, because the seller must
sell fast and the buyers who value them most are themselves short of funds.
\end{definition}

\begin{omfigure}
\begin{tikzpicture}[econ, every node/.style={font=\small}, box/.style={draw, rounded corners, align=center, text width=3.1cm,
    minimum height=1cm, fill=omDef!8}]
  \node[box] (a) at (0,2) {asset prices fall};
  \node[box] (b) at (5,2) {banks' equity falls};
  \node[box] (c) at (5,0) {leverage rises above its limit};
  \node[box] (d) at (0,0) {banks sell assets, often to each other};
  \draw[->, thick] (a) -- (b);
  \draw[->, thick] (b) -- (c);
  \draw[->, thick] (c) -- (d);
  \draw[->, thick] (d) -- (a);
\end{tikzpicture}

{\small The loop between leverage and forced sales. A bank that must restore its \omterm{def:b2:banks-and-crises:leverage}{leverage ratio} after a loss sells assets
or cuts lending; if many do so at once, prices fall further and the loss spreads to banks that held the same assets.}
\end{omfigure}

The loop works without any depositor running, and it works on lenders that are not banks: dealers, funds and others that
finance long assets with short borrowing. In 2007--08 several of the runs were on such short-term funding.

\section{2008 step by step}\label{sec:b2:banks-and-crises:2008}

\omimg{book4/northern-rock-queue-2007.jpg}{0.62\linewidth}{Customers outside a branch of Northern Rock, a British mortgage
lender, in Brighton on 14 September 2007, the day the Bank of England was authorised to support it. Photograph: Dominic
Alves, CC BY 2.0.}

\begin{record}[2008 step by step]
% ledger: b2:2008-bnp, b2:2008-northern-rock, b2:2008-bear, b2:2008-gse, b2:2008-lehman, b2:2008-aig-mmf, b2:2008-tarp
9 August 2007: BNP Paribas stops redemptions on three of its investment funds.
14 September 2007: the Bank of England is authorised to give liquidity support to Northern Rock, the fifth-largest British
mortgage lender. March 2008: the Federal Reserve finances the takeover of the investment bank Bear Stearns by JPMorgan
Chase, lending 29 billion dollars to a vehicle that holds 30 billion of its assets, and opens a lending facility for
securities dealers. 7 September 2008: the mortgage agencies
Fannie Mae and Freddie Mac are placed in government conservatorship. 15 September: Lehman Brothers files for bankruptcy.
16 September: the Federal Reserve lends up to 85 billion dollars to the insurer AIG, and the Reserve Primary Fund, a money
market fund, falls below one dollar a share on losses on Lehman's short-term debt. 19 September: the Treasury guarantees
money market funds. 3 October: a 700-billion-dollar programme to buy troubled assets is voted. 14 October: 250 billion
dollars of public capital is offered to banks, 125 billion of it to nine large institutions, and the debt of insured banks
is guaranteed.
\end{record}

Read with the model, much of 2008 was a run, though not by insured depositors: by holders of short claims on funds, dealers
and an insurer that had no \omterm{def:b2:banks-and-crises:deposit-insurance}{deposit insurance} and, before the crisis, no access to the central bank's lending. The public answers were the
ones of \cref{sec:b2:banks-and-crises:insurance}, extended to them: lending, guarantees, and new equity.

\section{Banks and government debt}\label{sec:b2:banks-and-crises:sovereign}

\begin{definition}[Bank--sovereign loop]\label{def:b2:banks-and-crises:doom-loop}
A \emph{bank--sovereign loop}\index{bank--sovereign loop} is a feedback in which losses at a country's banks raise its
government's debt, through support to the banks, while doubts about that debt lower the value of the government bonds the
banks hold.
\end{definition}

% ledger: b2:ft2018
Banks hold their own government's bonds as liquid assets, and governments stand behind their banks. When both are
strong, each supports the other. When one weakens, each drags the other down: support to the banks raises public debt and
its yield, the fall in bond prices weakens the banks, and the expected cost of supporting them rises again. Farhi and
Tirole (2018) built a model of this feedback, prompted by the euro-area crisis, in which the answers are domestic
support to the banks or support to the government from other countries.

\section{A century and a half of crises}\label{sec:b2:banks-and-crises:record}

\begin{definition}[Banking crisis]\label{def:b2:banks-and-crises:banking-crisis}
A \emph{banking crisis}\index{banking crisis} is an episode in which many of a country's banks fail or need public
support, with runs, large losses and a fall in lending; in the opening chart, a crisis is counted in the year it starts.
\end{definition}

\begin{omchart}{Real GDP per person in the five years after a business-cycle peak, log points relative to the peak year,
  averaged over 88 peaks with a banking crisis within two years and 203 other peaks; 18 advanced economies, 1870--2020,
  leaving out every peak whose six-year path touches a war year (1914--1919, 1939--1947).}{Jord\`a--Schularick--Taylor Macrohistory Database (rgdpmad, crisisJST); our computation}
\begin{tikzpicture}
\begin{axis}[omrecord, width=0.6\linewidth, height=0.4\linewidth, xmin=0, xmax=5, ymin=-4, ymax=6, xtick={0,...,5},
    xlabel={Years after the peak}, ylabel={Log points}, legend style={at={(1.03,0.5)}, anchor=west, font=\scriptsize}]
  \addplot[omThm, thick, mark=square*, mark size=1.3pt] table[x=h, y=normal] {figdata/bachelor-2/29-banks-and-crises/out/paths.csv};
  \addlegendentry{other peaks}
  \addplot[omDef, thick, mark=*, mark size=1.3pt] table[x=h, y=financial] {figdata/bachelor-2/29-banks-and-crises/out/paths.csv};
  \addlegendentry{peaks with a crisis}
\end{axis}
\end{tikzpicture}
\end{omchart}

% figdata: paths (financial -2.95, -2.70, -1.53, 0.01, 1.28; normal -2.18, -0.31, 1.62, 3.23, 4.89)
Recessions that come with a \omterm{def:b2:banks-and-crises:banking-crisis}{banking crisis} are a little deeper and much longer. One year after the peak, output per person is 2.95
log points below it with a crisis and 2.18 without; the difference is small. After a normal peak output is back above it
in the third year and 4.89 points above it in the fifth; after a crisis it is still 1.53 below in the third year and only
1.28 above in the fifth.

\begin{definition}[Credit boom]\label{def:b2:banks-and-crises:credit-boom}
A \emph{credit boom}\index{credit boom} is a period in which bank lending to households and firms grows much faster than
output, so that the ratio of credit to GDP rises fast.
\end{definition}

\begin{omchart}{Change in bank loans to the non-financial private sector as a share of GDP over five years, percentage
  points: the five years ending the year before the first year of a banking crisis (63 cases) and the same window before
  every other year (1,861); 18 advanced economies, 1870--2020, leaving out windows that touch a war year.}{Jord\`a--Schularick--Taylor Macrohistory Database (tloans, gdp,
  crisisJST); our computation}
\begin{tikzpicture}
\begin{axis}[omrecord, width=0.55\linewidth, height=0.4\linewidth, ybar, bar width=16pt, xmin=-0.6, xmax=1.6, ymin=0,
    ymax=15, xtick={0,1}, xticklabels={other windows, before a crisis}, ylabel={Percentage points},
    nodes near coords, nodes near coords style={font=\scriptsize}, legend style={at={(1.03,0.5)}, anchor=west,
    font=\scriptsize}]
  \addplot[fill=omDef!70, draw=omDef] table[x=group, y=mean] {figdata/bachelor-2/29-banks-and-crises/out/credit.csv};
  \addlegendentry{mean}
  \addplot[fill=omThm!50, draw=omThm] table[x=group, y=median] {figdata/bachelor-2/29-banks-and-crises/out/credit.csv};
  \addlegendentry{median}
\end{axis}
\end{tikzpicture}
\end{omchart}

% figdata: credit (before a crisis mean 13.35, median 10.55, n 63; other mean 4.79, median 4.41, n 1861)
% ledger: b2:st2012
In the five years before a crisis, loans rose by 13.35 points of GDP on average, against 4.79 in other five-year windows;
the medians are 10.55 and 4.41. Schularick and Taylor (2012), on 14 of these countries over 1870--2008, found credit growth
to be a powerful predictor of financial crises, and the output costs of crises to have remained large even as policy
responses became more aggressive. The chart compares averages: not every crisis followed a boom, and not every boom ended
in a crisis.

\section{Exercises}\label{sec:b2:banks-and-crises:exercises}

\begin{exercise}[$\star$]\label{exo:b2:banks-and-crises:1}
Read the opening chart. Which decade had the most crises, and which two came next? Which three had none?
\end{exercise}

\begin{exercise}[$\star$]\label{exo:b2:banks-and-crises:2}
From the chart of output paths: five years after the peak, how far above the peak is output per person after a normal peak,
and after a peak with a crisis? What is the gap?
\end{exercise}

\begin{exercise}[$\star$]\label{exo:b2:banks-and-crises:3}
A bank has assets of 100 and equity of 4. What share of its equity does a 2 per cent fall in asset values wipe out? And with
equity of 10?
\end{exercise}

\begin{exercise}[$\star$]\label{exo:b2:banks-and-crises:4}
Which events of 2007--08 in the record box were runs, and who was running from what?
\end{exercise}

\begin{exercise}[$\star\star$]\label{exo:b2:banks-and-crises:5}
Predict, then compute: with log utility, what is the optimal contract? Show that waiting is always a patient depositor's
\omterm{def:b2:static-games:best-response}{best response}.
\end{exercise}

\begin{exercise}[$\star\star$]\label{exo:b2:banks-and-crises:6}
With $\pi = 0.5$, $R = 2$ and $\gamma = 2$, find $c_1$, $c_2$ and the run threshold $f^*$.
\end{exercise}

\begin{exercise}[$\star\star$]\label{exo:b2:banks-and-crises:7}
Modify the listing to set $R = 1.5$ with $\gamma = 2$. Give $c_1$, $c_2$ and $f^*$, and say whether the bank is more or less
fragile than with $R = 2$.
\end{exercise}

\begin{exercise}[$\star\star$]\label{exo:b2:banks-and-crises:8}
Explain why suspension of convertibility removes the run when $\pi$ is known, and what goes wrong when it is not.
\end{exercise}

\begin{exercise}[$\star\star$]\label{exo:b2:banks-and-crises:9}
\textbf{(D)} With the Macrohistory data, count banking crises by decade, and compute the mean change in loans to GDP over the
five years ending the year before each crisis, leaving out windows that touch a war year. Compare with the charts.
\end{exercise}

\begin{exercise}[$\star\star\star$]\label{exo:b2:banks-and-crises:10}
Derive \eqref{eq:b2:banks-and-crises:contract}, and show that $c_1 > 1$ if and only if $\gamma > 1$.
\end{exercise}

\begin{exercise}[$\star\star\star$]\label{exo:b2:banks-and-crises:11}
Show that $w(f)$ is decreasing in $f$ if $c_1 > 1$, derive \eqref{eq:b2:banks-and-crises:threshold}, and show that
$f^* > \pi$ for the optimal contract.
\end{exercise}

\begin{exercise}[$\star\star\star$]\label{exo:b2:banks-and-crises:12}
The government promises every depositor who waits $c_2$, paid if needed by a tax on date-1 withdrawals. Show that no patient
depositor runs, and that the tax is never collected in equilibrium.
\end{exercise}

\section{Problem: The Brenland savings bank}\label{sec:b2:banks-and-crises:problem}

\begin{problem}[{Weekend problem --- The Brenland savings bank: depositors' types, the optimal deposit contract, the run
equilibrium, and insurance and suspension}]
\label{pb:b2:banks-and-crises:1}
A savings bank in a Brenland market town has depositors with one unit each. A quarter of them will need their money after
one season; the others can wait two. Money lent to farmers for two seasons doubles; recalled after one, it returns only what
was lent. Depositors have $u(c) = -1/c$.

\medskip
\noindent\textbf{Part I --- Depositors' types.}
\begin{enumerate}
  \item Give $\pi$, $R$ and $\gamma$.
  \item Without a bank, what does each type consume?
  \item What is a depositor's expected utility without a bank?
  \item Why can a depositor not buy insurance against being impatient from an insurer?
  \item Why can the bank not simply ask depositors their type?
\end{enumerate}

\medskip
\noindent\textbf{Part II --- The optimal deposit contract.}
\begin{enumerate}[resume]
  \item Write the bank's budget constraint.
  \item Write the first-order condition.
  \item Find $c_1$ and $c_2$.
  \item Check the budget constraint.
  \item Find the expected utility under the contract and compare it with Part I.
  \item Why would a patient depositor not pretend to be impatient?
\end{enumerate}

\medskip
\noindent\textbf{Part III --- The run equilibrium.}
\begin{enumerate}[resume]
  \item Write a waiting depositor's payoff if a share $f$ withdraws.
  \item What does a waiting depositor get if $f = 0.25$?
  \item If $f = 0.5$? If $f = 0.6$?
  \item Find the run threshold $f^*$.
  \item At what share of withdrawals is the bank empty?
  \item Show that everyone withdrawing is an equilibrium.
  \item What is lost in the run, compared with the good equilibrium?
\end{enumerate}

\medskip
\noindent\textbf{Part IV --- Insurance and suspension.}
\begin{enumerate}[resume]
  \item Under suspension, how much does the bank pay out at date 1?
  \item What does a waiting depositor then get, whatever others do?
  \item If a third of the depositors turned out to be impatient, what would suspension do to the last of them?
  \item Under a guarantee of $c_2$ to those who wait, what does the guarantee cost in equilibrium?
  \item Why might the guarantee change how the bank lends?
  \item If depositors were less risk-averse, with $\gamma = 1$, would the bank need any of this?
  \item State the result: the optimal first-period payout and the share of early withdrawals at which the bank fails.
\end{enumerate}
\end{problem}
